% ============================================================================
\chapter{Container Library}
% ============================================================================

\section{Dynamic Array (\texttt{Vector<T>})}

\texttt{Vector<T>} is a heap-growing array equivalent to \texttt{std::vector},
but with \texttt{IAllocator} injection. When no allocator is provided,
\texttt{operator new} / \texttt{operator delete} are used directly.

\subsubsection{Growth Strategy}

\begin{equation}
\text{newCapacity} =
\begin{cases}
  8 & \text{if current capacity} = 0 \\
  2 \times \text{current capacity} & \text{otherwise}
\end{cases}
\end{equation}

Doubling maintains amortised $O(1)$ push-back: a sequence of $n$
push-backs performs at most $2n$ element moves total.

\subsubsection{Construction and Destruction}

Elements are constructed in-place via placement-\texttt{new}:
\texttt{new (\&m\_data[m\_size]) T(value)}.
They are explicitly destroyed before deallocation:
\texttt{m\_data[i].\textasciitilde T()}.
This ensures correct behaviour for non-trivial types while remaining
compatible with the custom allocator interface.

\subsubsection{Move Semantics}

The move constructor and move assignment transfer ownership of the
underlying buffer in $O(1)$ by pointer swap, leaving the source in
a valid empty state.

\section{Hash Map (\texttt{HashMap<K,V>})}

The current \texttt{HashMap} is a \emph{linear-scan} map backed by a
\texttt{Vector<Pair>}. All operations are $O(n)$ in the number of entries.

\begin{remark}
The current implementation is intentionally simple: the expected
use-cases (editor panel registries, asset name$\to$ID tables) have at
most a few hundred entries where cache-friendly linear scan outperforms
hash table overhead. A true hash table (open addressing, Robin-Hood
probing) is planned for \texttt{v2.0} when scene complexity requires it.
\end{remark}

\section{String View (\texttt{StringView})}

\texttt{StringView} is a non-owning, read-only view into a null-terminated
string: a pointer + length pair. It performs no heap allocation. Comparison
is lexicographic and runs in $O(\min(|a|, |b|))$.

\section{Fixed String (\texttt{FixedString<N>})}

\texttt{FixedString<N>} stores at most $N-1$ characters in an inline
\texttt{char[N]} array plus a length field. There is no heap involvement.
It is used for entity names, asset paths, and any string that must be
embeddable inside a struct without pointer indirection.
